\section{The product expansion on a compact threefold}\label{sec:factor}

Let $X$ be a smooth projective Calabi--Yau threefold. Choose a
projective embedding of $X$ and use its degree on $\Gamma_X$.
Every nonempty effective cycle has positive degree. The integral
homology class is locally constant on the cycle space
\cite[\S2.1]{Pardon}. The Chow variety of cycles of each fixed degree
is compact, so its image in the discrete group $H_2(X;\ZZ)$ is finite.
Thus $\Gamma_X$ has finite degree sublevel sets, and the locus
$\mathscr C_\beta(X)$ of cycles of class $\beta$ is compact and
open and closed, including when different classes have the same
image modulo torsion.

Put $U_{X,\beta}=(X,1_{\mathscr C_\beta(X)})\in H$. Addition pulls the
characteristic function back to the sum over decompositions of $\beta$.
Under multiplication $D\mapsto dD$, the inverse image of
$\mathscr C_\beta(X)$ is the disjoint union of the
$\mathscr C_\gamma(X)$ over all effective $\gamma$ with $d\gamma=\beta$.
Pullback of its characteristic function counts each locus once.
This proves \eqref{eq:compat} for $\Z_X=\sum_\beta U_{X,\beta}Q^\beta$.

\begin{proof}[Proof of Theorem~\ref{thm:main}]
Use the symplectic form of the polarization. Proposition~3.3 and
Theorem~8.1 of Ionel and Parker \cite{IP} and the projective
comparison theorem of Siebert \cite[Theorem~0.1]{Siebert} give
\begin{equation}\label{eq:cluster}
 \log\GW(\Z_X)=
    \sum_{\beta\ne0}\sum_{g\ge0}c_{\beta,g}(X)\log E_g(Q^\beta,u).
\end{equation}
Here $\log\GW(\Z_X)$ is the connected positive-class series. The
integers in \eqref{eq:cluster} have finite genus support at each
$\beta$ by \cite[Definition~3.45 and Theorem~3.50]{DIW},
which gives finiteness for this elementary expansion itself.
Both the cluster expansion and the comparison of virtual classes
are indexed by $H_2(X;\ZZ)$. In particular \cite[(8.9)]{IP} sums
over every pair $(d,B)$ with $dB=A$ in integral homology. Siebert's
comparison fixes the integral class and the connected genus. The
semipositive description of the symplectic invariants is recalled
in \cite[Remark~1.1]{IP}.

Only classes of positive symplectic area can occur in the sum.
Suppose $\beta$ has positive ample degree and is ineffective for
the projective structure on $X$. Every $\gamma$ with $d\gamma=\beta$,
$d>1$, is then ineffective, since a positive multiple of an
effective cycle is effective. There are finitely many such pairs
because $d$ divides $\deg(\beta)$ and the roots for a fixed $d$,
if any, form a coset of the finite kernel of $d$ on $H_2(X;\ZZ)$.
Compare coefficients in
\eqref{eq:cluster}, inductively on the positive integral degree.
The left side vanishes at $\beta$, and all terms from proper
divisors already vanish. Thus
$\sum_g c_{\beta,g}(X) w^{g-1}=0$, using
$[q]\log E_g(q,u)=w^{g-1}$. Linear independence gives
$c_{\beta,g}(X)=0$ for every $g$.

The right side of \eqref{eq:main} is now well defined in
$H\llbracket\Gamma_X\rrbracket$. There are finitely many effective
classes of bounded degree, finitely many contributing genera per
class, and polynomial coefficients in each factor. Products and
inverses preserve the coproduct identity. By
\eqref{eq:operations}, $\rho_d A_g(Q^\beta)=A_g(Q^{d\beta})
=\psi_d A_g(Q^\beta)$. Both $\rho_d$ and $\psi_d$ are ring
endomorphisms, so the full product also satisfies the root sum
in \eqref{eq:compat}. Proposition~\ref{prop:evaluation} and
\eqref{eq:cluster} identify the disconnected GW evaluation of
the product with $\GW(\Z_X)$. Lemma~\ref{lem:reconstruction}
identifies the universal series. Uniqueness follows upon evaluation
from the unique elementary expansion, or by comparing the
multiplicity-one term at a first class of difference.
\end{proof}

The elementary-cluster evaluation in Corollary~\ref{cor:characterization}
is essential to the characterization. For example, fix $g\ne h$ and
replace $A_g(q)$ by $A_g(q)A_h(q^2)$, keeping every other factor.
The new series is compatible and has the same multiplicity-one
coefficient as $A_g$. In the compact product its exponent stays
$c_{\beta,g}$, and the exponent of $A_h(Q^\beta)$ becomes
$c_{\beta,h}-\sum_{2\gamma=\beta}c_{\gamma,g}$. Thus this change
preserves integer exponents in every class, although it changes the
degree-two evaluation of the genus-$g$ factor.

\subsection{Integrality from stable pairs}
There is a second way to see why the choice of action in
\eqref{eq:main} yields integer exponents. It uses the
Gromov--Witten/stable-pair correspondence for $H$ and the
form of the local hook series. We first obtain rational exponents
with finite genus support without using the numerical
cluster decomposition.

By Lemma~\ref{lem:reconstruction} the compatible $A_g(q)$ is a
product of rational powers of $f_h(q^d)$. Its degree-one evaluation
is $w^{g-1}$. Thus the entire factor at $d=1$ is $f_g(q)$, by the
independence of the powers $w^{h-1}$. All other factors have
$d\geq2$. At each multiplicity $m$ this gives
\begin{equation}\label{eq:triangular}
 a_{g,m}=x_{g,m}+P_{g,m}\bigl(\{x_{h,j}:1\le j<m\}\bigr)
\end{equation}
with $P_{g,m}$ a finite rational polynomial. Induction on $m$
inverts the substitution $x_{g,m}\mapsto a_{g,m}$. It is therefore
an automorphism of the polynomial algebra $H$. The identities
\eqref{eq:operations} show that it intertwines both the coproduct
and each $\rho_d$. Apply Lemma~\ref{lem:reconstruction} to the
inverse image of $\Z_X$, then apply the automorphism. This gives
a product in the $A_g(Q^\beta)$ with rational exponents
$n_{\beta,g}$ and finite genus support in every class.

Write $z$ for the stable-pair variable. Pardon's Theorem~1.7
\cite{Pardon}, extended by linearity over $\QQ$, identifies
$\PT(a_{g,m})$ with the expansion at $z=0$ of the rational
function $[q^m]E_g(q,u)$ under $e^{iu}=-z$. The identification
is unambiguous since expansion at $e^{iu}=1$ determines a rational
function. For a box of hook length $k$ one has
\[
 (2\sin(ku/2))^2=2-(-z)^k-(-z)^{-k}
                 =-(-z)^{-k}(1-(-z)^k)^2 .
\]
This is a unit in $\ZZ((z))$. Hence every coefficient of
$\PT(A_g(q))$ lies in $\ZZ((z))$, including when $g=0$ and
the hook factors are inverted. The compact stable-pair invariants
$\PT(U_{X,\beta})$ also lie in $\ZZ((z))$
\cite[\S3.5]{Pardon}.

Suppose all exponents in classes of smaller ample degree than
$\beta$ are known to be integral. Divide the stable-pair series by
the corresponding powers of the local factors. Their constant
terms are one, so the residual series still has coefficients in
$\ZZ((z))$. Its coefficient at $\beta$ is
\begin{equation}\label{eq:residual}
 \sum_g n_{\beta,g}\,z^{1-g}(1+z)^{2g-2},
\end{equation}
by the degree-one hook formula. In this finite sum the highest
genus with nonzero exponent has the unique lowest Laurent power,
and its local factor has leading coefficient one. The coefficient
of that power is therefore the exponent itself. It is integral.
Subtract its summand in \eqref{eq:residual} and continue in
descending genus. This proves integrality for every $\beta$.
Equation~\eqref{eq:cluster} identifies these integers as
$c_{\beta,g}(X)$.
